\begin{codebox}
\codeline{\textcolor{cmtgray}{\#\ ==============================}}
\codeline{\textcolor{cmtgray}{\#\ 5.\ CQ驱动的迭代流程}}
\codeline{\textcolor{cmtgray}{\#\ ==============================}}
\codeline{\textcolor{cmtgray}{\#\ 与领域专家访谈，收集CQ（10\textasciitilde{}30个为宜）}}
\codeline{\textcolor{cmtgray}{\#\ \ \ \(\rightarrow\)\ 按"基础/关联/校验/统计"分类，剔除超范围CQ}}
\codeline{\textcolor{cmtgray}{\#\ \ \ \(\rightarrow\)\ 反推类/属性/公理清单（见第3节）}}
\codeline{\textcolor{cmtgray}{\#\ \ \ \(\rightarrow\)\ 建模（七步法，见\ ontology-101-process.txt）}}
\codeline{\textcolor{cmtgray}{\#\ \ \ \(\rightarrow\)\ 每个CQ写验收查询，全部通过即完成一轮迭代}}
\codeline{\textcolor{cmtgray}{\#\ \ \ \(\rightarrow\)\ 新需求\ =\ 新CQ，进入下一轮迭代}}
\codeblank{}
\codeline{\textcolor{cmtgray}{\#\ 工程建议：}}
\codeline{\textcolor{cmtgray}{\#\ CQ清单应与本体一起进行版本管理（如\ cq/\ 目录下逐条编号），}}
\codeline{\textcolor{cmtgray}{\#\ 每次本体变更后自动回归执行全部验收查询（CI化，防止回归缺陷）}}
\end{codebox}
